% Section 1.8: Building Safe HAL Abstractions

\section{Implementing Hardware Abstraction Layer (HAL)}
\label{sec:hal-implementation}

A Peripheral Access Crate provides raw register access, but using it directly is error-prone and non-portable. The Hardware Abstraction Layer (HAL) transforms register-level operations into safe, ergonomic, and reusable APIs. This section guides you through building HAL abstractions for the EFR32MG24.

\subsection{From PAC to HAL: The Transformation}
\label{subsec:pac-to-hal}

The HAL sits between the PAC (register-level) and BSP (board-specific), providing:

\begin{itemize}
    \item Type-safe APIs that prevent misuse
    \item Zero-cost abstractions
    \item embedded-hal trait implementations
    \item Error handling and validation
    \item Consistent naming conventions
\end{itemize}

\subsection{HAL Design Principles}
\label{subsec:hal-design-principles}

\begin{designprinciple}{Safety First}
    Every HAL function must be memory-safe. Unsafe blocks are confined to the smallest possible scope and justified with comments.
\end{designprinciple}

\begin{designprinciple}{Type-State Pattern}
    Use Rust's type system to enforce state transitions at compile time. For example, a pin cannot be read when configured as output.
\end{designprinciple}

\begin{designprinciple}{Zero-Cost}
    HAL abstractions must compile to identical machine code as hand-written register operations.
\end{designprinciple}

\subsection{Implementing GPIO HAL}
\label{subsec:gpio-hal}

Let's build a GPIO HAL that demonstrates these principles:

\begin{minted}{rust}
pub struct Pin<PORT, PIN, MODE> {
    _mode: PhantomData<MODE>,
    _port: PhantomData<PORT>,
    _pin: PhantomData<PIN>,
}

impl<PORT, PIN> Pin<PORT, PIN, Input<Floating>> {
    pub fn into_output(self) -> Pin<PORT, PIN, Output<PushPull>> {
        critical_section::with(|_| {
            let gpio = unsafe { &*GPIO::ptr() };
            // Configure as output
            if PIN < 8 {
                gpio.model().modify(|r, w| {
                    w.bits((r.bits() & !(0xF << (PIN * 4))) | (0x4 << (PIN * 4)))
                });
            }
        });
        Pin { _mode: PhantomData, _port: PhantomData, _pin: PhantomData }
    }
}
\end{minted}

\subsection{The Role of embedded-hal Traits}
\label{subsec:embedded-hal-traits}

The \crate{embedded-hal} provides standardized traits:

\begin{minted}{rust}
// embedded-hal v1.0 traits
impl embedded_hal::digital::OutputPin for Pin<Output<PushPull>> {
    type Error = Infallible;
    
    fn set_low(&mut self) -> Result<(), Self::Error> {
        // Implementation
        Ok(())
    }
    
    fn set_high(&mut self) -> Result<(), Self::Error> {
        // Implementation
        Ok(())
    }
}
\end{minted}

\subsection{Clock Management HAL}
\label{subsec:clocks-hal}

Clock configuration requires careful sequencing:

\begin{minted}{rust}
pub struct Clocks {
    _private: (), // Prevent direct instantiation
}

impl Clocks {
    pub fn new() -> Self {
        // Critical: Enable HFXO first
        let cmu = unsafe { &*CMU::ptr() };
        cmu.hfxoctrl.write(|w| w.mode().xtal());
        while cmu.hfxostatus.read().hfxordy().is_not_ready() {}
        
        Self { _private: () }
    }
}
\end{minted}

\subsection{Error Handling in HAL}
\label{subsec:hal-errors}

HALs should provide meaningful error types:

\begin{minted}{rust}
#[derive(Debug, Clone, Copy)]
pub enum HalError {
    Timeout,
    InvalidConfiguration,
    DeviceBusy,
    BufferOverflow,
}

type HalResult<T> = Result<T, HalError>;
\end{minted}

\begin{example}
Example: UART with timeout
\end{example}

\begin{minted}{rust}
impl Uart {
    pub fn transmit(&mut self, data: &[u8]) -> HalResult<()> {
        for byte in data {
            // Check if TX buffer has space
            if self.tx_buffer_full() {
                return Err(HalError::BufferOverflow);
            }
            self.send_byte(*byte)?;
            
            // Wait with timeout
            let mut attempts = 1000;
            while self.transmit_in_progress() && attempts > 0 {
                attempts -= 1;
            }
            if attempts == 0 {
                return Err(HalError::Timeout);
            }
        }
        Ok(())
    }
}
\end{minted}

\begin{keytakeaways}
    \item HAL provides safe abstractions over raw PAC
    \item Type-state pattern prevents misuse at compile time
    \item embedded-hal traits enable ecosystem compatibility
    \item Zero-cost abstractions maintain performance
    \item Explicit error handling improves debuggability
\end{keytakeaways}

\begin{exercises}
    \item\textbf{Easy}: Implement a safe GPIO input API using type-state pattern.
    
    \item\textbf{Medium}: Create an I2C HAL that enforces bus arbitration rules at compile time.
    
    \item\textbf{Hard}: Design a HAL for DMA that prevents buffer overflow through type system guarantees.
\end{exercises}
